\def\TheFile{ch01_intro.tex}

% quote text
\begin{savequote}[15cm]
  \raggedleft
  \sffamily
  Introducing oneself properly is always hard.
  \qauthor{Anonymous}
\end{savequote}
\chapter{Introduction}

Writing documentation is often considered a chore. But actually reading student reports is even worse.
Certainly if the student is to wordy, sloppy, repeats every other section and so on.

The Casus Belli in this case is that I have been examiner of a lot,
not to say most students in the informatics courses at Fontys
Hogeschool Venlo. There I have to read 22 reports, of 14 students that
I coach and 8 others where I am the examiner. That incentivesed me to
write down some advice. Here you have it.

\section{Do not bore us to death!}

\begin{wrapfigure}{r}{.25\textwidth}
  \includegraphics[width=.25\textwidth]{images/boreme.png}
\end{wrapfigure}
The worst thing that can happen to me, the poor person that has to
read your whole report from front page until the last page before the
appendix, is that I get the impression that some one is trying to hum
me asleep.  I have no clue what or who caused this, but many students
have the habit to explain in the first paragraph of the chapter what
they are going to tell in that chapter. Why? Why Waste my time and your time? Do you think that improves your grade?
Because I have to read the whole damn thing anyway, why take the excitement of
reading away. Would you read a book or watch a film if the first
paragraph of each chapters gives everything away?

\section{Things to avoid}

To keep the reader awake, and more importantly: interested, and apreciative of your work you should:
\paragraph{Stay DRY} DO NOT repeat stuff.
\paragraph{Good titles} Think of good chapter and section titles. We read and use the table of content. If the chapter and section titles are good, they help explain the structure of the report, without any extra boring help.
\paragraph{Be brief} You do not get paid per written word, nor are we paid per word \textbf{read}.
\paragraph{Use a storyline} Both in your report and in your presentation. If you invent a persona anyway, use him or her as the protagonist to tell the store, and use him/her to explain stuff.
The story need not necessarily be true to the actual chronology of the time spent in your bachelor project, but should be logical story, in which you take the reader along to explain your reasoning, the decisions you made and why, etc.

You protagonist may not be useful in all the parts, but might be very handy to explain the problem, the assignment and how he or she can you the fruits of your labor.

Start with the \textbf{company} and every detail of the company that is relevant to the project.
Then continue with the \textbf{context} of the problem, exactly as much as you need to explain the next logical thing: the \textbf{problem} that you have been asked to solve. Followed by the \textbf{assignment}.

If in any paragraph that you wrote you or a reviewer notices that you have to explain things
that could have been explained in an earlier part, do not repeat any of that earlier part,
not even by saying 'as you can/have read in ..', but instead revise the earlier part so the context, the problem, the assignment etc can be understood in one flow.

As an example, If you need to explain what a typical customer will do with the product, then we expect that you have explaind product and customer in the company description or context. 

You should use the reader as if he/she is a single pass compiler. What has not been defined before cannot be used here. Forward references (like in as you will see in ...) are \textbf{not allowed}. It is a waste of words, in particular when the structure, and thus the table of content is aby good.
That will keep the reader in the flow, because he does not have to page forward or backward, and if the story is short enough, the reader will be able to pull through without being bored to death.

Also: By NOT repeating stuff, you have no chance of inconsistencies.

\section{TLDR;}

In the remainder of this document you will see some tips and tricks to use when you write your report in \LaTeX,
but the above and some things in the use of graphics also apply when you use \textbf{word} or some other text processing application. The quality of your report should not suffer from your choice of tools.
So read those chapters as well.

\section{Use a better technology}
One of the standards for documentation in open source and hence in
Linux land is \LaTeX, a text processing package. \LaTeX\ is available
for free and available with all Linux distributions and installed in
the Lab.

\TeX\ is the machinery of \LaTeX\ and was defined in the \TeX\ book
\cite{texbook} and implemented by
Prof. Donald Knuth. \LaTeX\ is a (nowadays HUGE) set of macros built
on top of that. \LaTeX\ in its initial form is described by Leslie
Lamport in \cite{latexbook}. If you like your book thick, try the
\LaTeX\ companion \cite{latexcompanion}.

The web is also a very good source of \LaTeX\ documentation. A good
starting point is \url{http://en.wikibooks.org/wiki/LaTeX}, useful for
beginners and pros alike.

This is a simple multi part document. It's purpose is to show how easy it is
to create a multi part document, one that, for instance, can be worked on 
simultaneously by several authors. Note that most of the settings for
this document are set in the file \texttt{mydefs.tex}. 
Look in that file too.

%% This document consists of the following files
%% \begin{enumerate}[noitemsep,topsep=0pt,parsep=0pt,partopsep=0pt]
%% \item \texttt{main.tex}
%% \item \texttt{motivation.tex}
%% \item \texttt{graphics.tex}
%% \item \texttt{codelisting.tex}
%% \item \texttt{poser.tex}
%% \item \texttt{hello.c}
%% \item \texttt{Makefile} and
%% \item \texttt{mydefs.tex}
%% \item \texttt{servlet3.png}, the pie chart
%% \item \texttt{Diagram1.pdf}, the UML diagram
%% \end{enumerate} 

You are kindly advised to keep your lab logs in simple text
files. These can be turned into latex files easily,
which can be used to produce a nice looking report.
\clearpage 
\section{Some hints to start with}
Sometimes things do not work out the way you think.
\LaTeX\ interpretes some character codes in it's own way.
Things like dollar signs or even underscore are special.
\LaTeX\ source are littered with accolades or \textit{curly braces} if that's
the way you call them. They are special too. So here is some advice: 

Do not use \define{funny file names}. That is: stick to ASCII filenames without spaces or even underscores. 
These will bring only you into trouble. If you want to keep things portable, 
don't use camel case (like in JavaClassNames) either, because
some OS-es do not distinguish between upper and lower case. You may of
course brake this rule if the files are program things like
Java source files.

%% \subsection{Hints for informatics (use version control)}
%% In software projects, versioning is important. \LaTeX\ and \textsc{SubVersion}
%% work nicely together here.
%% By using the \LaTeX\ package \texttt{svn-multi} you can have svn+
%% latex insert meta information about your files in the produced output,
%% for instance in the page footers, as in this document.

%% If you add the codes below at the top of each .tex file, these
%% codes will be expanded/updated by svn \textit{on checkin}. 
%% % do not expand keywords on  the example file
%% \lstinputlisting{svnkeywords.head}
%% You must also tell
%% subversion to expand these keywords for those files with for each file
%% the command:\\
%% {\lstset{language=sh,numbers=none,morekeywords={svn,propset,keywords,filename},keywordstyle={\color{blue}}}
%% \begin{lstlisting}[frame=single]
%% svn propset svn:keywords "Id Author File Date LastChangeDate
%%   Revision HeadURL Header"  filename
%% \end{lstlisting}
%% }
%% To keep these version codes up to date, first check if your \LaTeX\  files compile,
%% then check them in and do your final \LaTeX\  run. 

%% To make the version codes of the files appear on the bottom line, have
%% a look at the mydefs.tex file where the fancy headers and footers are defined.
%% In group work you may also find it comfortable to create a file named
%% \texttt{myauthors.tex} with contents like below, which is then
%% input into the main or mydefs at the proper place and can translate
%% student numbers from svn and you peerweb account into a humanly
%% readable authors name.

%% \lstinputlisting[frame=single]{myauthors.tex}

%%% Local Variables: 
%%% mode: latex
%%% TeX-master: "main"
%%% End: 
